We fix Grothendieck universes $\U\in\V$. Categories of shapes and spaces are
$\U$-small; categories of coefficients may be $\V$-large. A construction
requiring an additional universe says so explicitly. We put
$BX=|NX|$, $e=\{*\}$, and
$\Hot=\Cat_{\U}[\Winfty^{-1}]$, where $\Winfty$ consists of the
functors whose nerves are weak homotopy equivalences. We write $\W_n$
for the functors inducing isomorphisms on $\pi_i$ for $0\leq i\leq n$,
at all base points; the assertion on $\pi_0$ is bijectivity.

Coefficients have the manuscript's presheaf variance. Thus the enhanced
example is $\D_{\mathcal C}(A)=\Ho\Fun((NA)^{\op},\mathcal C)$.
In formulas about spaces we identify a groupoid with its opposite by
inverting paths. For $u:A\to B$, the comma category $A/b$ has objects
$(a,u(a)\to b)$, and $b\backslash A$ has objects $(a,b\to u(a))$.
A covariantly final functor has every $b\backslash A$ nonempty and
connected. A sieve is closed under sources, and a cosieve under targets.

In this paper a \emph{fundamental localizer} satisfies identities,
two-out-of-three and closure under retracts; contains $A\to e$ when
$A$ has a terminal object; and satisfies the \emph{relative} Theorem A:
if $u:A\to B$ is a morphism over $S$ and every $A/s\to B/s$ belongs
to $\W$, then $u\in\W$. In Chapter XVI this last axiom is literally
Loc~(3~bis), whereas Loc~(3) is its absolute special case. This distinction
is material below. An object is $\W$-aspherical if its specified map to
$e$ is in $\W$.

Every numbered assertion carries one of the requested statuses.
\emph{Open} means unresolved by this work, not a claim of literature-wide
openness. \emph{Proved} records an argument given here and makes no
priority claim. \emph{Known} records a checked source, with the precise
version used. Earlier AI-produced companion files are audited inputs,
not primary mathematical authorities. The checked CATHOT C.1--C.4 are
used as supplied in the prompt. All original arguments and all limitations
in this document are attributable only to the credited AI system.
